\SourceLeaf{407}%
\[
\text{V.}\quad .\ .\ .\ .\ .\ .\ .\qquad b=90^0,\quad A\gtrless90^0
\]
\par
1) \(c=90^0\). Dann ist \(A\) der Pol von \(c\), also gilt \(b=90^0\) und \(a=A\). Die Gleichung
\[
\cos a=\cos b.\cos c+\sin b.\sin c.\cos A
\]
ist unmittelbar ersichtlich.

2) \(c<90^0\). Hier verlängert man \(c\) über \(B\) hinaus bis zur Länge \(90^0\). Aus der Betrachtung des Dreiecks folgt dann
\[
\cos a=\cos(90^0-c).\cos A+\sin(90^0-c).\sin A.\cos90^0
\]
oder
\[
\cos a=\sin c.\cos A
\]
Somit ist die Formel auch in diesem Fall richtig.

3) \(c>90^0\). Hier schneidet man von \(c\) den Bogen \(90^0\) im Punkt \(R\) ab. Aus der Betrachtung des Dreiecks \(BRC\) folgt dann
\[
\cos a=\cos(c-90^0).\cos A+\sin(c-90^0).\sin A.\cos90^0
\]
oder
\[
\cos a=\sin c.\cos A
\]
wie im vorigen Fall.

Die Fälle, in denen \(b<90^0\) oder \(b>90^0\) und zugleich \(c=90^0\) gilt, sind im Wesentlichen dieselben wie die beiden vorigen.

Die Aufzählung aller möglichen Fälle liefert somit den Beweis.
\begin{center}
\begin{tabular}{ll}
Beide Seiten kleiner als \(90^0\)&aus I\\
Beide Seiten größer als \(90^0\)&aus II\\
Eine größer, die andere kleiner als \(90^0\)&aus III\\
Beide \(=90^0\)&aus IV und V.\,1.\\
Eine \(=90^0\), die andere kleiner&aus IV und V.\,2.\\
Eine \(=90^0\), die andere größer&aus IV und V.\,3.\\
\end{tabular}
\end{center}
\hspace*{\parindent}Wir können daher allgemein annehmen, dass gilt
\[
\text{(A)}\quad .\ .\ .\ .\ .\quad\cos a=\cos b.\cos c+\sin b.\sin c.\cos A
\]
und ebenso
